In this section, we develop a supply-side model which, together with our demand estimates, facilitates counterfactual simulations to evaluate the welfare and profit impact of the monitoring program. We begin by comparing the status quo regime to a scenario in which monitoring is not available. This comparison establishes a baseline quantification for the value that the monitoring program provides. We then examine how this value depends on the firm's pricing strategy, competitive pricing responses, and potential data portability regulations.

\subsection{A Model of Firm Pricing with the Monitoring Program}\label{subsec:pricing}

As we discussed in Sections \ref{subsec:bg_pricing} and \ref{sec:demand_model}, special programs like monitoring are typically priced as multiplicative factors to a baseline price menu $p_{t=0}(y|x)$ across coverage options $y$ and observable classes $x$. In keeping with this, we model firms' pricing strategies through a vector of multiplicative adjustment factors that scale the baseline price menus observed in our data. We term these adjustment factors as pricing \textit{levers} and denote them by $\boldsymbol{\kappa}$ to distinguish them from other (fixed) pricing factors. This approach allows us to abstract away from regulatory and actuarial details that are largely unrelated to monitoring, such as the risk-rating associated with driver age or vehicle types.
Accordingly, we denote the observed baseline price menu for a representative individual by $p$, suppressing the notation for $x$ and $y$.
Similarly, we denote draws of renewal factors $r_{t}(m|s,x,\nu)$ as defined in \Cref{eq:renewal-factor-score} by $r(m|s)$. The observed baseline renewal price menu is thus $p\cdot r(m|s)$. Because monitoring only affects renewal factors in the first period,
we omit the period index $t$.

The goal of our exercise is to study how the monitoring firm can affect equilibrium outcomes by differentiating prices across the monitored and unmonitored pools, and among monitored consumers after the monitoring period concludes. To this effect, we consider pricing strategies parameterized by a vector of four levers $\boldsymbol{\kappa} = \{\kappa_{0s}, \kappa_{0d},\kappa_{1s}, \kappa_{1d}\}$. The first lever $\kappa_{0s}$ multiplies all prices, setting a new baseline for all drivers. The second lever $\kappa_{0d}$ multiplies the initial price offered to monitored drivers, generating an opt-in discount. The remaining levers $\kappa_{1s}$ and $\kappa_{1d}$ modify the renewal price that monitored drivers are offered depending on the realization of their monitoring scores. Together, the pricing levers $\boldsymbol{\kappa}$ generate the following progression of prices for prospective consumers.

\paragraph{Initial pricing for unmonitored drivers ($t = 0$)}
In their first period with the firm, a consumer who opts out of monitoring receives the modified baseline price. That is, if the consumer were offered a price $p$ under the status quo, then their initial period price would become $\kappa_{0s} \cdot p$ under the pricing strategy $\boldsymbol{\kappa}$.

\paragraph{Renewal pricing for unmonitored drivers ($t \geq 1$)}
At renewal, unmonitored consumers are offered the baseline progression of renewal prices.\footnote{At renewal, the firm has no incentive to modify its baseline pricing for unmonitored customers: it gains no additional information about these drivers, nor does it have reasons to nudge them toward a particular plan choice as it did with monitoring in the first period.} That is, at their first renewal, a consumer who did not enter monitoring would draw a renewal factor $r(0|0)$ according to the distribution described in \Cref{eq:renewal-factor-score}. Absent a claim, their price for the following period would be $\kappa_{0s} \cdot p \cdot r(0|0)$. If a claim were to occur, then their renewal price would further increase by the claim factor $r_{c}$ defined in \Cref{eq:renewal-rate-overview}.

\paragraph{Initial pricing for monitored drivers ($t=0$)}

A consumer who opts in to monitoring receives a \textit{one-time} discount in the first period. That is, if the consumer were offered a price $p$ under the status quo, then their price for the monitoring period would become $\kappa_{0s} \cdot \kappa_{0d} \cdot p$ under the pricing strategy $\boldsymbol{\kappa}$.

\paragraph{Renewal pricing for monitored drivers ($t \geq 1$)}

The monitoring score generated in the initial period provides the firm with a posterior update of each monitored driver's private risk type. To leverage this information, the firm applies a different lever for drivers whose scores outperform or underperform the firm's expectations.  For drivers who outperform expectations (e.g., they received a monitored renewal factor lower than the unmonitored one $r(1|s) \leq r(0|0)$ in the data), the firm applies lever $\kappa_{1d}$, whereas for drivers who underperform, the firm applies  $\kappa_{1s}$.
We model the renewal factor in each case as a mixture between $r(1|s)$ and $r(0|0)$ as follows:
\begin{align}\label{eq:renewal-factor-kappa}
r^{\ast}(m=1|s) =
\begin{cases}
\kappa_{1d} \cdot r(1|s) + (1- \kappa_{1d}) \cdot r(0|0) & \text{if }r(1|s)\leq r(0|0)\\
\kappa_{1s} \cdot r(1|s) + (1- \kappa_{1s}) \cdot r(0|0) & \text{if }r(1|s)>r(0|0).
\end{cases}
\end{align}
Here, $\kappa_{1d}$ can be interpreted as a  \textit{rent-sharing} parameter for monitored consumers who prove to be safer drivers.
For instance, if a consumer were to receive a 30\% discount at renewal after monitoring in our data, $\kappa_{1d}=0$ would reduce their discount to 0 so that their renewal price would be equal to that of a driver who did not enter monitoring.  Because this consumer is known to be less costly to insure than the average unmonitored driver, the increase in renewal price would be accrued directly to firm profits.
If $\kappa_{1d}=1.5$, their discount would instead increase by a half to 45\%, so that a larger share of the cost savings revealed by monitoring were given back to the consumer. Similarly, $\kappa_{1s}$ controls the degree to which monitored consumers who prove to be costlier to insure may see higher prices---a \textit{risk surcharge} parameter.

Note that because the opt-in discount $\kappa_{0d}$ is only applied in the first period, the baseline price at renewal is the same for monitored and unmonitored drivers. As such, a consumer who completes their monitoring period with score $s$ and no accidents would expect a renewal price of $\kappa_{0s} \cdot p \cdot r^{\ast}(1|s)$ in the subsequent period.

\paragraph{Firm Profits and Variable Costs}

To choose $\boldsymbol{\kappa}$ optimally, the firm maximizes its aggregate expected profits across all consumers in the market. For each consumer, it internalizes their risk and preference types, integrating over possible realizations of claims, renewal price shocks, and scores to anticipate their claims costs, initial-period plan choice, renewal price changes and possible attrition or resorting into different coverage levels. We assume that the firm does not experience time-discounting, but that its profit planning horizon matches consumers' welfare horizon of three periods in our baseline model.

The marginal profit from a consumer in a given period is given by the expected premium under each plan, less the expected cost of insuring them given the plan's coverage, multiplied by the probability of that consumer choosing the plan. We assume that the cost of insuring a consumer in a given period is determined by the sum of their expected claims in that period (subject to the deductible and limits of their chosen plan) and an additive \textit{loading cost} as defined in \Cref{eq:pricing_def}, reflecting overhead as well as expenses to administer the plan and potential claims. For monitored drivers and during the initial period only, the firm benefits from reduced risk but incurs a \textit{variable cost of monitoring}. Based on our interviews with the firm's management team, this variable cost is about \$35.

\paragraph{Discussion} The four pricing levers in our model allow the firm to exercise an ``invest-and-harvest'' strategy that generates ex-ante ambiguous welfare effects for different types of drivers. The renewal levers $\kappa_{1s}$ and $\kappa_{1d}$ determine how the firm prices the monitoring {data}. A lower rent-sharing lever $\kappa_{1d}$ or higher surcharge lever $\kappa_{1s}$ transfers less of the informational surplus generated by the monitoring technology to consumers, representing a higher level of extraction from monitored drivers---``harvesting'' the data---but it also discourages drivers from opting in to monitoring in the first place. These considerations may affect drivers with higher or lower risk asymmetrically. The profit-maximizing degree of advantageous selection is therefore an empirical question: it depends, among other factors, on how the additional risk information revealed by monitoring correlates with risk aversion, price sensitivity, and monitoring disutility.

Conditional on the renewal levers, $\kappa_{0s}$ and $\kappa_{0d}$ function as ``investment'' levers to encourage initial enrollment in monitoring. A higher $\kappa_{0s}$ and a lower $\kappa_{0d}$ encourage more drivers to opt in, generating more data and potential for rent extraction in renewal periods. The trade-off is lower inframarginal profit in the initial period and the risk of losing monitoring-averse consumers to competitors. As such, the optimal levels of $\boldsymbol{\kappa}$ depend on the competitive landscape as well.

\subsection{Equilibrium Game}\label{subsec:equi_concept}

The monitoring firm's four pricing levers interact strategically with its competitors' pricing decisions. To capture these equilibrium dynamics, we consider a Bertrand pricing game in which prices are set simultaneously by our firm and a \textit{focal} competitor.

Recall from \Cref{sec:demand_model} that our competitive landscape consists of the three largest firms in the market---the monitoring firm and two competitors whose prices enter our demand estimation in renewal choices. Among these competitors, one offers the lowest average price and the other offers the median among all competitors in our dataset. For our main results, we assume that the focal competitor is the low-price firm. The median-price competitor is treated as passive, giving consumers a realistic outside option whose price remains unaffected by monitoring. To match observed market shares, this passive firm is modeled as a composite residual firm that absorbs all market shares unserved by the two active players.
\Cref{subsec:model_sensitivity} examines robustness to alternative market and competitor definitions.

\paragraph{Competitor Pricing} As discussed in \Cref{subsec:data_monitoring}, we assume that competitors do not offer monitoring on their own. However, the focal competitor can adjust its pricing in response to the monitoring program in two ways. First, the competitor can adjust its baseline pricing via a baseline lever $\kappa^c_0$ akin to $\kappa_{0s}$. Second, if monitoring data becomes fully portable so that the competitor is able to observe and interpret a driver's monitoring score in the same way that the monitoring firm can, then  the competitor can also counter the monitoring firm's renewal offer to monitored drivers using a competing set of levers $\kappa^c_{1s}$ and $\kappa^c_{1d}$, parameterized similarly to their counterparts:
\begin{align*}
r^{c\ast}(m=1|s) & =
\begin{cases}
\kappa^c_{1d} \cdot r(1|s) + (1- \kappa^c_{1d}) \cdot r(0|0) & \text{if }r(1|s)\leq r(0|0)\\
\kappa^c_{1s} \cdot r(1|s) + (1- \kappa^c_{1s}) \cdot r(0|0) & \text{if }r(1|s)>r(0|0).
\end{cases}
\end{align*}

Under this model, a driver considering the focal competitor without monitoring would face a competitor price of $p^c \cdot \kappa^c_{0s}$. A monitored driver considering renewal with score $s$ would face a competitor price of $p^c \cdot \kappa^c_{0s} \cdot r^{c\ast}(1|s)$.

\paragraph{Counterfactual Scenarios and Regulations} We consider five counterfactual scenarios. ``Status Quo'' reflects the most recent baseline and monitoring pricing regimes for all firms, which covers the majority of our data after the introduction of monitoring.
The ``No Monitoring'' scenario removes the monitoring program altogether. As we discuss in Appendix \ref{sec:app_firm_pricing}, the introduction of monitoring did not alter baseline prices. We thus assume that both the monitoring firm and the focal competitor would have continued to set their baseline levers $\kappa_{0s}$ and $\kappa^c_{0s}$ to be 1 in this counterfactual.\footnote{In practice, the pricing of the monitoring program, including its interaction with baseline prices for unmonitored drivers, was coarsened and simplified relative to a profit-maximizing price schedule in order to appeal to consumers and regulators \parencite{ben2023privacy}. In Appendix \ref{sec:app_firm_pricing}, we show that the discount schedule for monitored drivers was coarse and lumpy (\Cref{fig:score-discount-by-regime}), while the baseline prices for unmonitored drivers were also totally unaffected by the introduction of monitoring (\Cref{fig:app_event_study}b).} Comparing these two scenarios allows us to evaluate the impact of the monitoring program.

Next, we consider counterfactual scenarios in which the monitoring firm sets profit-maximizing levers $\boldsymbol{\kappa}$. First, a ``Partial'' equilibrium scenario holds fixed the focal competitor's pricing. We then allow the latter to respond in the ``Full'' equilibrium scenario. However, as monitoring data is proprietary to the firm, the focal competitor can only set $\kappa^c_{0s}$.

A recurring policy proposal in the context of monitoring is to mandate data portability.\footnote{For example, the \href{https://home.treasury.gov/system/files/311/181031_EU\%20US\%20Big\%20Data\%20\%20Issue\%20Paper\_Publication.pdf}{EU-U.S. Insurance Dialogue Project's 2018 Big Data issue paper} (with NAIC participation) identified data portability as an emerging consumer issue for U.S. state insurance supervisors. The Federal Trade Commission's
\href{https://www.ftc.gov/system/files/documents/public_events/1568699/transcript-data-portability-workshop-final.pdf}{``Data To Go: An FTC Workshop on Data Portability.'' (2020)}
further examined data portability's potential benefits and risks for competition and switching.
} Just as accident records must be observable to all insurance firms, a data portability law would require firms to share any monitoring data that they collect with the whole industry (at least upon a drivers' request). The reasoning behind this proposal is that data portability would increase competition and mitigate asymmetric information in the insurance market: if competitors are able to discern less costly drivers as well as the monitoring firm can, then they can undercut the monitoring firm's renewal prices, limiting its ability to extract rents. However, if competition at renewal eliminates the firm's ability to profit from the information generated by the monitoring program, it may not be incentivized to invest in the program in the first place. Moreover, while competitors would benefit from additional data on driver risk, the pricing of such information can further heighten reclassification risk associated with monitoring, thereby reducing opt-in. Nonetheless, monitoring can remain attractive under a data portability mandate for two reasons. First, the moral hazard effect and reduced risk during the monitoring period is unaffected by the mandate. Second, large switching frictions give the monitoring firm a competitive advantage in retaining safe drivers at renewal, even if its information advantage is removed.

To assess this policy, we consider a ``Data Port.'' counterfactual. This scenario equalizes firms' information sets at renewal and allows them to compete using the full set of post-monitoring price levers. In order to simulate realistic regulatory constraints, we restrict the domain of feasible price levers for the firms in two ways: the ex-post risk surcharge cannot exceed what we observe in the data ($\kappa_{1s}, \kappa^{c}_{1s} \leq 100\%$), and the rent sharing factor must be non-negative ($\kappa_{1d}, \kappa^{c}_{1d} \geq 0$). These constraints reflect the regulatory principle that monitoring programs should primarily reward safe drivers through premium reductions rather than penalize risky drivers \parencite{ben2023privacy}.\footnote{For example, New York's regulatory framework explicitly requires usage-based insurance programs to be ``discount only'' so as not to ``negatively impact policyholders.''} Taking the constraints imposed by regulation one step further, we consider a final ``Disc. Floor'' scenario that also restricts rent sharing be no less than the level observed in the data ($\kappa_{1d}, \kappa^{c}_{1d} \geq 100\%$), essentially enforcing a discount floor.

\paragraph{Loading Cost Calibration}
We do not observe the loading costs for either the monitoring firm or its competitors. As such, we apply the following calibration procedure to approximate them. Modeling each firm's loading cost as a uniform per-consumer additive scalar, we maintain the standard assumption that the baseline prices observed for each firm maximize the firm's expected profits given its loading cost. Since the pricing of monitoring had no bearing on baseline prices in reality, we calibrate loading costs using unmonitored drivers only. Specifically, we perform a one-dimensional grid search for each firm, looking for a value that rationalizes each of the observed baseline prices ($\kappa_{0s} = \kappa^{c}_{0s} = \kappa^{c,\text{passive}}_{0s} = 1$) as mutual best responses in the ``No Monitoring'' scenario.

\paragraph{Calibrating Firm Brand Effects}
Our dataset covers the customers of the monitoring firm only. To conduct market-level counterfactual simulations, we assume that these consumers are representative of the full market and compute outcomes on a per-consumer, per-year basis. Specifically, we assume that
(1) the distribution of risk and preference types in our sample matches that of the full market; and (2) in the first period, consumers choose one the three firms in addition to coverage levels and monitoring participation according to our model in \Cref{eq:dynamic_decision_model}, with the modification that their flow utilities also include firm-level (``brand'') fixed effects. To calibrate these brand effects, we normalize the monitoring firm's brand effect to zero, and conduct a two-dimensional grid-search for the brand effects of the two competitors that best match the status quo market shares. Since the passive competitor is a composite residual firm, its market share and the brand effect will naturally be large.

\paragraph{Grid Search and Multiple Equilibria}
Due to complex interactions between consumer costs, demand and the multi-plan choice menu offered by each firm, the payoff space for each firm may not be concave with respect to its pricing parameters. As such, it is possible for multiple equilibria (or no pure strategy equilibria) to exist. To identify equilibrium outcomes for each of our counterfactual regimes, we perform an exhaustive grid search over plausible price parameter tuples.
Whenever there are multiple equilibria, we select the one that maximizes the monitoring firm's profit.

\subsection{Counterfactual Simulation Results}

For each of the five counterfactual scenarios in \Cref{tab:ctf_main}, we present results in three categories: (1) surplus and surplus division into consumer welfare, focal firm and competitor profits; (2) quantity measures such as insurance coverage, market share, and monitoring adoption; and (3) pricing levers, given by the $\boldsymbol{\kappa}$ and $\boldsymbol{\kappa^c}$, subject to regulatory caps.

All welfare and profit measures are reported in annualized, per-capita dollars, computed by scaling the per-period total quantities to match the yearly time frame and then averaging across all consumers in the market (\textit{not} weighted by the likelihood of coming to the monitoring firm). Because consumer welfare is inherently nonlinear, we report certainty-equivalent changes relative to the ``no monitoring'' scenario. ``Total surplus'' equals the sum of the changes in consumer welfare and industry profits.

As part of the quantity measures, we report the expected amount of insurance coverage (in thousands of dollars) associated with consumers' plan choices, regardless of firm. We also report the market share of the monitoring program overall and separately among new customers (``Initial-Period'') and existing customers (``Renewal''), all in percentage terms.

The pricing levers for the monitoring firm ($\boldsymbol{\kappa}$) and the competitor ($\boldsymbol{\kappa^c}$) are reported directly. For instance, a baseline lever $\kappa_{0s}$ of $1.09$ corresponds to a 9 percentage point surcharge over the observed lever of $1.00$. Similarly, an initial-period monitoring lever  $\kappa_{0d}$ of 0.23 corresponds to a 77\% discount relative to baseline pricing, and a 73 percentage point discount over the observed lever of 4\%.

\inputifexists{\exhibitpath/tab_7}

\paragraph{Assessing the Monitoring Program}
The first two columns of \Cref{tab:ctf_main} illustrate the welfare impact generated by the monitoring program, as it is implemented in our data. Although the monitoring firm only captures an 11.85\% market share, the average consumer in the entire insurance market gains \$7.43 per year in certainty equivalent from the availability of monitoring. The monitoring firm gains \$5.53 while the competitor firms lose \$3.18 per-capita per-year. Summing across all parties, the combined total surplus improves by \$9.79 per-capita per-year due to the introduction of the monitoring program.

\paragraph{Optimal Pricing of the Monitoring Program}
In the ``Partial'' equilibrium regime, profit-maximizing pricing amplifies the ``invest-and-harvest'' dynamic for the monitoring firm. The firm raises baseline prices by 9\% while offering a 77\% one-time discount to encourage opt-in, regardless of the ex-post monitoring outcome. However, the firm shares little rent with safer drivers (only 0.05) compared to the 1.00 benchmark observed in the data, exploiting its information advantage and consumers' high switching costs. This is notable because baseline prices are already higher than observed in the data. For riskier consumers, however, the surcharge also falls from 1.00 in the data to 0.58, indicating that the firm still expects positive profits from most of them despite higher risk.

On net, firm profits increase by 8.0\% relative to the status quo, eroding some of the consumer welfare gains from monitoring. Nonetheless, total surplus rises as monitoring adoption more than doubles to 6.91\% of the overall market, and the average consumer chooses a higher coverage level.

In the ``Full'' equilibrium scenario, the focal competitor responds to the monitoring firm's optimal pricing by lowering baseline prices by just 1\%. The monitoring firm reacts by scaling back ex-ante ``investment'': lowering both the baseline surcharge and the upfront opt-in discount. However, ex-post rent extraction is exacerbated: safe drivers receive 0\% rent sharing, while risky drivers face a percentage point higher surcharge.

\paragraph{Data Portability} When consumer data portability is enforced, the competitor can in principle selectively ``poach'' consumers after monitoring. However, our results show that the only equilibrium is a corner solution in which the competitor poaches as little as possible: no discount for safe drivers and maximum surcharge for risky drivers. This weakly increases the competitor's pricing for monitored drivers, which depresses monitoring opt-in ex-ante by raising reclassification risk in the event of a bad monitoring score. The monitoring firm thus responds by slightly increasing its upfront discount and lowering its ex-post surcharge.

Why is poaching unsustainable in equilibrium? The payoff matrix reveals that, as soon as the competitor sets an interior value for $\kappa^c_{1s}$ and $\kappa^c_{1d}$ to attract monitored drivers, the monitoring firm can always profitably counter by lowering its baseline surcharge ($\kappa_{0s}$) or raising its upfront discount ($\kappa_{0d}$). Both strategies steal market share from the competitor while having opposite effects on monitoring opt-in and inframarginal profits. This allows the firm to fine-tune its response to, and effectively counter, any poaching.

In equilibrium, data portability affects consumer welfare primarily by increasing reclassification risk, rather than by intensifying competition. But this does not mean that poaching has no effect. Our final counterfactual scenario adds a mandated discount floor requiring that rent-sharing with safe drivers by either firm be no less than the observed level ($\kappa_{1d}, \kappa^{c}_{1d} \geq 1.00$x). This effectively forces the competitor to poach safe drivers with a high $\kappa^c_{1d}>1$. The results show that the monitoring firm counters lower ex-post rent extraction by lowering ex-ante ``investments,'' raising its initial-period surcharge and lowering discounts, ultimately reducing welfare and profits.

Taken together, our results show that ex-post price controls---capping the surcharge for risky drivers, setting safe driver discount floors, or data portability mandates---can harm firms' dynamic incentives to invest in monitoring while raising reclassification risk. In our case, the welfare benefits of monitoring through accident reduction and information generation are so large that these harms overwhelm the benefits of increased competition. As such, these regulations wind up harming consumers.

\subsection{Sensitivity and Discussion}\label{subsec:model_sensitivity}

In this section, we test the robustness of our results to several key modeling assumptions. We focus on alternatives to our specification of (1) the firm's variable monitoring cost, (2) how forward-looking consumers and firms are,
(3) regulatory constraints on monitoring pricing, and (4) pre-existing market structure and the focal competing firm.

Overall, our findings of substantial welfare benefits from the monitoring program and of harms from a prospective data portability regulation remain qualitatively identical and quantitatively robust to our main specification. Most quantities remain highly stable across specifications, while industry profits and firms' pricing levers vary more noticeably in magnitude, albeit only in one alternative specification each.

\paragraph{Marginal Cost of Monitoring}
Tables \ref{tab:ctf_robust_low_cost} and \ref{tab:ctf_robust_high_cost} present scenarios in which the marginal cost of monitoring is 50\% higher or lower than our baseline of \$35. Our results remain qualitatively similar to the baseline.

\paragraph{Time Horizons}
In Tables \ref{tab:ctf_robust_2p} and \ref{tab:ctf_robust_4p}, we show that our results are qualitatively identical when we adopt the two-period and four-period look-ahead models. Quantitatively, both scenarios see lower consumer welfare gains compared to the baseline, but for different reasons. A shorter horizon reduces the future benefits of monitoring, diminishing the firm's investment and monitoring opt-in rates. A longer horizon does the opposite, but the first-period-only claim reduction effect is also spread out to more periods, ultimately reducing the average per-period gain. In addition, although the data sharing equilibrium still features the corner solution of minimal poaching from the competitor, the monitoring firm now counters primarily by raising its baseline surcharge while scaling back upfront discounts. This triggers a higher baseline surcharge by the competitor, further eroding consumer welfare.

\paragraph{Ex-post Price Regulations}
Our baseline simulations assume that the rent-sharing factors ($\kappa_{1d}, \kappa^{c}_{1d}$) are non-negative, and that the risk-surcharge factors ($\kappa_{1s}, \kappa^{c}_{1s}$) are no higher than the level observed in the data. Table \ref{tab:ctf_robust_unconstraint} shows that, when we relax these constraints to +/- 200\%, our results remain very stable. In particular, the data portability regime sees another corner solution in which the focal competitor does not poach by charging the maximum possible price for both safe and risky drivers after monitoring. This reinforces the idea that the monitoring firm's pricing levers are flexible enough to effectively counter any ex-post competitor poaching in equilibrium, preserving the firm's ability to offer and profit from monitoring.\footnote{In our testing, this result holds even if the competitor does not commit to future prices, and if we further widen the constraints.}

\paragraph{Pre-existing Market Structure}
The monitoring firm only has 12\% market share in our data. Our baseline results also pit it against the lowest-pricing competitor. What if its baseline market position were enhanced? Or if our model allowed it to directly compete against other firms?

Table \ref{tab:ctf_robust_minmed_only} presents results in which the market consists of the monitoring firm, and the minimum- and median-pricing competitors only. By replacing the composite residual competitor with the medium-pricing competitor, the monitoring firm's market share increases to 27\%. While the results remain qualitatively similar to the baseline, both the benefit of monitoring and the firm's pricing of it are more pronounced.

Table \ref{tab:ctf_robust_med} simulates a scenario in which we switch the focal competitor with the passive composite competitor that has 66\% market share. Facing a more dominant firm, the monitoring firm scales back its upfront discount but increases ex-post rent-sharing to retain safe drivers. This reduces the monitoring opt-in rate relative to the baseline in our main specification. The surplus impact of monitoring thus decreases. The welfare impact, however, increases whenever the competitor is allowed to lower its prices in equilibria and decreases otherwise.
